% !PRINTED 128
\weberpage{128}
%<<CONT
die Indices von $Q$ und den mit $Q$ conjugirten Theilern von $P$ in Bezug auf $Q$. Daher ist $h_1 = 1$ und

\[
p^{\alpha-\beta} = h_1 + h_2 + h_3 \ldots % [sic: Weber prints no plus sign before the ellipsis; the preceding decomposition requires one.]
\]

Die Summanden $h_1, h_2, h_3, \ldots$ können hier als Theiler des Grades von $Q$ nur Potenzen von $p$ sein, und da $h_1 = 1$ ist; so folgt, dass mindestens $p$ dieser Summanden $= 1$ sein müssen, da sonst ihre Summe nicht durch $p$ theilbar sein könnte. Das heisst aber nichts Anderes, als dass wenigstens $p - 1$ der conjugirten Gruppen $b^{-1}Qb$, wo $b$ nicht in $Q$ enthalten ist, durch $Q$ theilbar und also gleich $Q$ sein müssen. Wenn aber

\[
b^{-1}Qb = Q
\]

ist, so ist auch für jeden Exponenten $s$

\[
b^{-s}Qb^s = Q.
\]

Ist $b^r$ die niedrigste Potenz von $b$, die in $Q$ enthalten ist, so muss, da der Grad von $b$ eine Potenz von $p$ ist, auch $r$ eine Potenz von $p$ sein, und

\[
b^{\frac{r}{p}} = c
\]

ist ein Element von $P$, das selbst nicht in $Q$ vorkommt, dessen $p^{\text{te}}$ Potenz aber in $Q$ enthalten ist, und das zugleich der Bedingung

\[
cQc^{-1} = Q
\]

genügt. Daraus folgt, dass die Elemente

\[
R = Q + Qc + Qc^2 + \cdots + Qc^{p-1}
\]

eine Gruppe vom Grade $p^{\beta+1}$ bilden, von der $Q$ ein Normaltheiler ist. Wir haben also den Satz:

III. \gesperrt{Ist $P$ eine Gruppe vom Grade $p^\alpha$ und $Q$ ein Theiler von $P$ vom Grade $p^\beta$ ($\beta < \alpha$), so giebt es einen Theiler $R$ von $P$ vom Grade $p^{\beta+1}$, von dem $Q$ Normaltheiler ist.}

Nimmt man in diesem Satze $\beta = \alpha - 1$ an, so fällt $R$ mit $P$ zusammen, und es ergiebt sich, dass $Q$ ein Normaltheiler von $P$ ist. Wendet man denselben Satz auf die Gruppe $Q$ an u. s. f., so erhält man:

IV. \gesperrt{Jede Gruppe, deren Grad eine Primzahlpotenz ist, ist metacyklisch (\S.\,8).}
